%! AP Statistics — Unit 2, Lesson 2.1 — Student Guided Notes
\documentclass[10pt]{article}
\usepackage{apstats-article}
\usepackage{apstats-boxes}

%=============================================================================
\begin{document}

\pageheader{Unit 2, Lesson 2.1}{Guided Notes}
\namedateperiod

\vspace{0.12in}

% ── OBJECTIVE ────────────────────────────────────────────────────────────────
\begin{objectivebox}
\small By the end of this lesson, I will be able to\dots\\[4pt]
\writeline\\[1pt]
\writeline\\[1pt]
\writeline
\end{objectivebox}

\vspace{0.1in}

% ── VOCABULARY ───────────────────────────────────────────────────────────────
\begin{vocabbox}
\termblanklong{Two-variable data}
\termblanklong{Association}
\termblanklong{Apparent association}
\termblanklong{Causation}
\termblanklong{Confounding variable}
\termblanklong{Explanatory variable}
\termblanklong{Response variable}
\end{vocabbox}

\vspace{0.1in}

% ── HOOK ─────────────────────────────────────────────────────────────────────
\begin{hookbox}
\small\textbf{Headlines:}
\begin{itemize}[leftmargin=1.3em, itemsep=2pt]
  \item \textit{``Students who eat breakfast score higher on tests.''}
  \item \textit{``Shoe size and reading ability are strongly associated in children.''}
\end{itemize}

\vspace{6pt}
\textbf{Discussion --- write your thoughts:}
\begin{enumerate}[leftmargin=1.5em, itemsep=2pt]
  \item Could the breakfast--grades relationship be real? Could it be random? What would make you sure?\\
        \writeline\\[1pt]\writeline

  \item Shoe size and reading ability are correlated in children. But shoe size does not teach children to read.
        What might actually explain this relationship?\\
        \writeline\\[1pt]\writeline
\end{enumerate}
\end{hookbox}

\vspace{0.1in}

% ── FROM ONE VARIABLE TO TWO ──────────────────────────────────────────────────
\begin{notesbox}{From One-Variable to Two-Variable Questions}
\small
\renewcommand{\arraystretch}{1.5}
\begin{tabularx}{\linewidth}{@{} >{\columncolor{sky}\bfseries\color{navy}}p{0.36\linewidth}
                                    X @{}}
\textbf{Unit 1 question (one variable)} & \textbf{Unit 2 question (two variables)} \\
\hline
What is the typical GPA of students at our school?
  & Is there a relationship between sleep and GPA? \\
How many hours of sleep do students get?
  & Do students who play sports sleep more or less than those who don't? \\
\end{tabularx}

\vspace{8pt}
\textbf{The shift:} Unit 2 asks whether two variables \blank{3cm} --- whether knowing
one variable tells us something about the \blank{3cm}.
\end{notesbox}

\vspace{0.1in}

% ── APPARENT VS. REAL ASSOCIATION ────────────────────────────────────────────
\begin{notesbox}{Apparent Associations May Be Random (VAR-1.D.1)}
\small
\begin{multicols}{2}

\textbf{Key idea:}\\
A pattern visible in a \blank{2cm} may disappear when
we collect \blank{2cm} data.
An apparent association is \textbf{not} guaranteed to
be a real relationship in the \blank{2cm}.

\vspace{6pt}
\textbf{Example:} Suppose in a class of 10 students,
the 3 who drink coffee all have higher GPAs.
Does that prove coffee helps GPA?\\[4pt]
\writeline\\[1pt]\writeline

\columnbreak

\textbf{What makes an association ``real'' vs.\ ``random''?}
\begin{itemize}[itemsep=3pt]
  \item Larger \blank{2.5cm} $\to$ less likely to be random
  \item Stronger \blank{2.5cm} $\to$ less likely to be random
  \item Consistent across \blank{2cm} groups $\to$ more credible
\end{itemize}

\vspace{6pt}
\textbf{Bottom line:} Statistics gives us tools to decide.
We'll build those tools in Units 5--9.
\end{multicols}
\end{notesbox}

\vspace{0.1in}

% ── ASSOCIATION vs. CAUSATION ─────────────────────────────────────────────────
\begin{notesbox}{Association vs.\ Causation}
\small
\begin{multicols}{2}

\textbf{Three explanations for an observed association:}
\begin{enumerate}[leftmargin=1.4em, itemsep=4pt]
  \item $X$ \blank{2.5cm} $Y$
  \item $Y$ \blank{2.5cm} $X$
  \item A third variable \blank{2cm} both $X$ and $Y$
        (this is called a \blank{3cm} variable)
\end{enumerate}

\columnbreak

\textbf{Classic example --- fire trucks:}\\
\textit{Cities that send more fire trucks to a fire have
more property damage.}\\[6pt]
Does this mean fire trucks \emph{cause} damage?\\
\writeline\\[1pt]\writeline\\[1pt]
What is the confounding variable?\\
\writeline

\end{multicols}
\end{notesbox}

\vspace{0.1in}

% ── GUIDED PRACTICE ──────────────────────────────────────────────────────────
\begin{practicebox}
\small\textbf{Research Question Analysis}

\vspace{5pt}
\textit{Question: ``Do students who spend more time on social media report lower
levels of life satisfaction?''}

\vspace{6pt}
\begin{tabularx}{\linewidth}{@{} l X @{}}
\textbf{Two variables:}           & \blank{9cm} \\[6pt]
\textbf{Types (cat/quant):}       & \blank{9cm} \\[6pt]
\textbf{Could the pattern be random?} & \blank{9cm} \\[6pt]
\textbf{One confounding variable:} & \blank{9cm} \\[6pt]
\textbf{Proves causation?}         & \blank{9cm} \\[6pt]
\end{tabularx}
\end{practicebox}

\vspace{0.1in}

% ── UNIT 2 ROAD MAP ──────────────────────────────────────────────────────────
\begin{spiralbox}
\small
\begin{multicols}{2}

\textbf{Where we go from here:}
\begin{itemize}[itemsep=3pt]
  \item Both categorical $\to$ \textbf{Lessons 2.2--2.3} (two-way tables, bar charts)
  \item Both quantitative $\to$ \textbf{Lessons 2.4--2.9} (scatterplots, regression)
  \item \textit{Establishing causation} $\to$ \textbf{Unit 3} (experimental design)
\end{itemize}

\columnbreak

\textbf{Big Idea --- write it in your own words:}\\
\textit{Why can't we always trust an apparent association?}\\[2pt]
\writeline\\[1pt]\writeline\\[1pt]\writeline

\end{multicols}
\end{spiralbox}

\vspace{0.1in}

% ── EXTRA NOTES ──────────────────────────────────────────────────────────────
\begin{notesbox}{Additional Notes}
\vspace{0pt}
\writeline\\\writeline\\\writeline\\\writeline\\\writeline\\\writeline
\end{notesbox}

\end{document}
